\newpage\handout
{Drill Problems 11}
{\textsc{Scholastic Aptitude Test (SAT)}}
{\href{https://creativecommons.org/licenses/by-nc-sa/4.0/}{CC BY-NC-SA 4.0 license}}
{Author: \BookAuthor}{Release: \generatedOn}
{SAT: Drill Problems (11)}
\input{drill00-disclaimer.tex}


\begin{enumerate}
  \item \textbf{Solving for an Expression} (10 points)\\
  Given that $2(x+4)=18$, what is the value of $x+6$?
  \begin{enumerate}[label=\Alph*.]
    \item $-4$
    \item $5$
    \item $9$
    \item $11$
  \end{enumerate}
  \begin{subanswer}
    % your answer here
  \end{subanswer}

  \item \textbf{Inequality Solution} (10 points)\\
  $$3x+2<-5(x+6)$$\\
  Which of the following values is a solution to the given inequality?
  \begin{enumerate}[label=\Alph*.]
    \item $-5$
    \item $-4$
    \item $3$
    \item $4$
  \end{enumerate}
  \begin{subanswer}
    % your answer here
  \end{subanswer}

  \newpage

  \item \textbf{Mean from Bar Graph} (10 points)\\
  \insertimage{0.40}{images/cus18}{reference attached}
  The bar graph shows how many hours the air conditioner was on from Monday to Friday in a given week. What is the mean number of hours the air conditioner was on during this period?
  \begin{enumerate}[label=\Alph*.]
    \item $8$
    \item $10$
    \item $12$
    \item $14$
  \end{enumerate}
  \begin{subanswer}
    % your answer here
  \end{subanswer}

  \item \textbf{Quadratic Expansion Constant} (10 points)\\
  The expression $-2(x+3)^2+6$ is equivalent to $a x^2+b x+c$, where $a<0$ and $b<0$. What is the value of $c$?
  \begin{subanswer}
    % your answer here
  \end{subanswer}

  \item \textbf{Slope of Perpendicular Line} (10 points)\\
  What is the value of the slope of a line perpendicular to the line $5x-6y+30=0$?
  \begin{subanswer}
    % your answer here
  \end{subanswer}

  \newpage

  \item \textbf{Sum of Quadratic Roots} (10 points)\\
  $$3x^2+9x-27=0$$\\
  If $m$ and $n$ are solutions to the given equation, what is the value of $m+n$?
  \begin{subanswer}
    % your answer here
  \end{subanswer}

  \item \textbf{Point on a Circle} (10 points)\\
  The graph of $(x-3)^2+y^2+8y=84$ is a circle in the $xy$-plane. Which of the following coordinates lie on the circle?
  \begin{enumerate}[label=\Alph*.]
    \item $(1,7)$
    \item $(-2,5)$
    \item $(-3,4)$
    \item $(3,-6)$
  \end{enumerate}
  \begin{subanswer}
    % your answer here
  \end{subanswer}

  \item \textbf{Stockpile Percentage Remaining} (10 points)\\
  At the start of the winter, the US stockpile of corn was 5 billion bushels. Two months later, $p\%$ of the stockpile had been used. Which expression best represents the number, in billions, of bushels of corn remaining in the US stockpile at the end of those two months?
  \begin{enumerate}[label=\Alph*.]
    \item $2(5)\left(\frac{100-p}{100}\right)$
    \item $5\left(\frac{100-p}{100}\right)$
    \item $\left(\frac{p}{100}\right)$
    \item $2(5)(100-p)$
  \end{enumerate}
  \begin{subanswer}
    % your answer here
  \end{subanswer}

  \item \textbf{Profit Function Interpretation} (10 points)\\
  The equation $P(y)=37(1.1)^y$ gives the estimated annual profit, in thousands of dollars, at a pottery studio, where $y$ is the number of years since the studio moved to a new location. Which of the following is the best interpretation of the number 37 in this context?
  \begin{enumerate}[label=\Alph*.]
    \item The estimated annual profit, in thousands of dollars, when the studio moved to the new location
    \item The increase in the estimated annual profit, in thousands of dollars, each year
    \item The number of years since the studio moved to the new location
    \item The percent increase in the estimated annual profit each year
  \end{enumerate}
  \begin{subanswer}
    % your answer here
  \end{subanswer}

  \item \textbf{Isosceles Triangle Angle} (10 points)\\
  Triangle $JKL$ is an isosceles triangle. The measure of angle $J$ is $30^\circ$. If $JL$ is the longest side of the triangle, what is the measure of angle $K$?
  \begin{enumerate}[label=\Alph*.]
    \item $30$
    \item $75$
    \item $120$
    \item $150$
  \end{enumerate}
  \begin{subanswer}
    % your answer here
  \end{subanswer}

  \item \textbf{Function Value from Rule} (10 points)\\
  If $f(2x)=9x-7$, what is the value of $f(6)$?
  \begin{enumerate}[label=\Alph*.]
    \item $11$
    \item $20$
    \item $47$
    \item $101$
  \end{enumerate}
  \begin{subanswer}
    % your answer here
  \end{subanswer}

  \item \textbf{System of Equations Solutions} (10 points)\\
  \[
    \begin{gathered}
    3x-4y=16 \\
    -6x=-8y+32
    \end{gathered}
  \]
  How many solutions does the given system of equations have?
  \begin{enumerate}[label=\Alph*.]
    \item Exactly one
    \item Exactly two
    \item Infinitely many
    \item Zero
  \end{enumerate}
  \begin{subanswer}
    % your answer here
  \end{subanswer}

  \item \textbf{Solving for Product} (10 points)\\
  If $\frac{5}{x}=\frac{2}{y}$ and $\frac{4}{x}-\frac{2}{y}=-\frac{1}{5}$, what is the value of $xy$?
  \begin{subanswer}
    % your answer here
  \end{subanswer}

  \item \textbf{Radical and Exponent Equation} (10 points)\\
  In the equation $\frac{1}{2} \sqrt{a}=2 \sqrt[3]{b}$, both $a$ and $b$ are positive real numbers. If $a=8$ and the equation is written in the form $a^x=b$, what is the value of $x$?
  \begin{subanswer}
    % your answer here
  \end{subanswer}

  \item \textbf{Exponential Function from Points} (10 points)\\
  The function $g$ is defined by $g(x)=a^x-b$, where $a$ and $b$ are constants. In the $xy$-plane, the graph of $g(x)$ passes through the points $(0,-6)$ and $(1,2)$. What is the value of $a$?
  \begin{enumerate}[label=\Alph*.]
    \item $6$
    \item $7$
    \item $8$
    \item $9$
  \end{enumerate}
  \begin{subanswer}
    % your answer here
  \end{subanswer}

  \item \textbf{Quadratic Solution with Radical} (10 points)\\
  If $1+\frac{a \sqrt{2}}{2}$ is a solution to the equation $2x^2-4x-7=0$ and $a>0$, what is the value of $a$?
  \begin{subanswer}
    % your answer here
  \end{subanswer}

  \item \textbf{Transformation of Quadratic Function} (10 points)\\
  If $f(x)=2(x-3)^2+8$ is transformed into $g(x)=2(x-5)^2+5$, which of the following describes the transformation?
  \begin{enumerate}[label=\Alph*.]
    \item The $x$-coordinate moves to the right 2 units and the $y$ coordinate moves 3 units down.
    \item The $x$-coordinate moves to the left 2 units and the $y$-coordinate moves 3 units down.
    \item The $x$-coordinate moves to the right 2 units and the $y$ coordinate moves 3 units up.
    \item The $x$-coordinate moves to the left 2 units and the $y$-coordinate moves 3 units up.
  \end{enumerate}
  \begin{subanswer}
    % your answer here
  \end{subanswer}

  \newpage

  \item \textbf{Unit Conversion: Density} (10 points)\\
  A sample of dried yellow pine has a density of 420 kilograms per cubic meter. To the nearest tenth, what is the density, in pounds per cubic foot, of this sample? (Use 1 kilogram = 2.2 pounds and 1 meter $=3.3$ feet)
  \begin{enumerate}[label=\Alph*.]
    \item $5.3$
    \item $25.7$
    \item $127.3$
    \item $280.0$
  \end{enumerate}
  \begin{subanswer}
    % your answer here
  \end{subanswer}

  \item \textbf{Arc Length in Circle} (10 points)\\
  \insertimage{0.26}{images/cus19}{reference attached}
  In the given circle, the radius is 6 centimeters and angle $AOB$ measures $\frac{7}{12}\pi$ radians. What is the value of the arc length $AB$, in centimeters?
  \begin{enumerate}[label=\Alph*.]
    \item $\frac{7}{2}\pi$
    \item $\frac{21}{2}\pi$
    \item $12\pi$
    \item $36\pi$
  \end{enumerate}
  \begin{subanswer}
    % your answer here
  \end{subanswer}

  \newpage

  \item \textbf{Right Triangle Altitude Length} (10 points)\\
  \insertimage{0.50}{images/cus20}{reference attached}
  In the image shown, right triangle $LMN$ has a right angle at $M$. Altitude $PM$ is shown and $\angle NPM$ is a right angle. $LP=13$ and $PM=10$. Which of the following are closest to the length of $NP$?
  \begin{enumerate}[label=\Alph*.]
    \item $3.0$
    \item $6.2$
    \item $7.1$
    \item $7.7$
  \end{enumerate}
  \begin{subanswer}
    % your answer here
  \end{subanswer}

  \item \textbf{Quadratic from Table: Find $b$} (10 points)\\
  $f(x)=24x^2+bx+c$\\
  For the given function $f$, $b$ and $c$ are constants, and the graph of $y=f(x)$ in the $xy$-plane passes through the points in the table below. What is the value of $b$?
  \begin{table}[h!]
  \centering
  \renewcommand{\arraystretch}{1.3}
  \setlength{\tabcolsep}{16pt}
  \begin{tabular}{|c|c|}
  \hline $x$ & $f(x)$ \\
  \hline $-\frac{15}{8}$ & 0 \\
  \hline 0 & $c$ \\
  \hline $\frac{11}{4}$ & 0 \\
  \hline
  \end{tabular}
  \end{table}
  \begin{subanswer}
    % your answer here
  \end{subanswer}

  \newpage

  \item \textbf{Polynomial Factors from Table} (10 points)\\
  \begin{table}[h!]
  \centering
  \renewcommand{\arraystretch}{1.3}
  \setlength{\tabcolsep}{16pt}
  \begin{tabular}{|c|c|}
  \hline $x$ & $g(x)$ \\
  \hline $-7$ & $10$ \\
  \hline $-4$ & $0$ \\
  \hline $0$ & $5$ \\
  \hline $3$ & $1$ \\
  \hline $5$ & $-6$ \\
  \hline $11$ & $0$ \\
  \hline
  \end{tabular}
  \end{table}
  Select values for the polynomial function $g(x)$ are shown in the table. Based on the values in the table, which of the following must be factors of $g(x)$?
  \begin{enumerate}[label=\Alph*.]
    \item $(x-5)$
    \item $(x+5)$
    \item $(x-4)$ and $(x+11)$
    \item $(x+4)$ and $(x-11)$
  \end{enumerate}
  \begin{subanswer}
    % your answer here
  \end{subanswer}

\end{enumerate}

